% Ad Familiares 9.24 (ad-familiares-09-24)
% Source: https://raw.githubusercontent.com/PerseusDL/canonical-latinLit/master/data/phi0474/phi056/phi0474.phi056.perseus-lat1.xml
% Fetched by scripts/fetch_latin.py

% Dateline (Perseus): Scr. Romae ante mcd. m. Febr. a. 711 (43) .

\ciceroLetterOpener{CICERO PAETO S. D.}

\ciceroSection{1} Rufum istum, amicum tuum, de quo iterum iam ad me scribis, adiuvarem, quantum possem, etiam si ab eo laesus essem, cum te tanto opere viderem eius causa laborare ; cum vero et ex tuis litteris et ex illius ad me missis intellegam et iudicem magnae curae ei salutem meam fuisse, non possum ei non amicus esse, neque solum tua commendatione, quae apud me, ut debet, valet plurimum, sed etiam voluntate ac iudicio meo. volo enim te scire, mi Paete, initium mihi suspicionis et cautionis et diligentiae fuisse litteras tuas quibus litteris congruentes fuerunt aliae postea multorum. nam et Aquini et Fabrateriae consilia sunt inita de me, quae te video inaudisse, et, quasi divinarent quam iis molestus essem futurus, nihil aliud egerunt nisi me ut opprimerent. quod ego non suspicans incautior fuissem, nisi a te admonitus essem. quam ob rem iste tuus amicus apud me commendatione non eget. utinam ea fortuna rei p. sit, ut ille me unum gratissimum possit cognoscere! sed haec hactenus.

\ciceroSection{2} te ad cenas itare desisse moleste fero ; magna enim te delectatione et voluptate privasti ; deinde etiam vereor (licet enim verum dicere) ne nescio quid illud, quod solebas, dediscas et obliviscare, cenulas facere. nam si tum, cum habebas quos imitarere, non multum proficiebas, quid nunc te facturum putem? Spurinna quidem, cum ei rem demonstrassem et vitam tuam superiorem exposuissem, magnum periculum summae rei p. demonstrabat, nisi ad superiorem consuetudinem tum, cum Favonius flaret, revertisses ; hoc tempore ferri posse, si forte tu frigus ferre non posses.

\ciceroSection{3} sed me hercule, mi Paete, extra iocum moneo te, quod pertinere ad beate vivendum arbitror, ut cum viris bonis, iucundis, amantibus tui vivas. nihil est aptius vitae, nihil ad beate vivendum accommodatius. nec id ad voluptatem refero sed ad communitatem vitae atque victus remissionemque animorum, quae maxime sermone efficitur familiari, qui est in conviviis dulcissimus, ut sapientius nostri quam Graeci ; illi sumpo/sia aut su/ndeipna , id est compotationes aut concenationes, nos 'convivia,' quod tum maxime simul vivitur. vides ut te philosophando revocare coner ad cenas. cura ut valeas ; id foris cenitando facillime consequere.

\ciceroSection{4} sed cave, si me amas, existimes me, quod iocosius scribam, abiecisse curam rei p. sic tibi, mi Paete, persuade, me dies et noctes nihil aliud agere, nihil curare, nisi ut mei cives salvi liberique sint. nullum locum praetermitto monendi, agendi, providendi ; hoc denique animo sum ut, si in hac cura atque administratione vita mihi ponenda sit, praeclare actum mecum putem. etiam atque etiam vale.
